% --- Archon physics formalization source begin ---
% archon:physics
% archon:covers PhyXMiniProblems/problem_phyx_mini_0502.lean
% archon:source-report reports/phyx_mini/problem_phyx_mini_0502.source.json

\chapter{Physics problem phyx\_mini\_0502}
\label{ch:PhyXMiniProblems_problem_phyx_mini_0502}

\paragraph{Problem source.}
\#\# Physical scenario

Figure shows the probability density for an electron that has passed through an experimental apparatus.

\#\# Question

What is the probability that the electron will land in a \textbackslash{}( 0.010\textbackslash{}text\{-mm\} \textbackslash{})-wide strip at \textbackslash{}( x = 0.000\textbackslash{}text\{ mm\} \textbackslash{})?

\#\# Answer choices

- A: A: \textbackslash{}( 3.0 \textbackslash{}times 10\textasciicircum{}\{-4\} \textbackslash{})

- B: B: \textbackslash{}( 3.0 \textbackslash{}times 10\textasciicircum{}\{-3\} \textbackslash{})

- C: C: \textbackslash{}( 5.0 \textbackslash{}times 10\textasciicircum{}\{-3\} \textbackslash{})

- D: D: \textbackslash{}( 5.0 \textbackslash{}times 10\textasciicircum{}\{-4\} \textbackslash{})

\#\# Figure caption (auxiliary; use the image as primary evidence)

The image is a plot labeled \textbackslash{}(P(x) = |\textbackslash{}psi(x)|\textasciicircum{}2 \textbackslash{}, (\textbackslash{}text\{mm\}\textasciicircum{}\{-1\})\textbackslash{}). The x-axis is labeled as \textbackslash{}(x \textbackslash{}, (\textbackslash{}text\{mm\})\textbackslash{}) and ranges from \textbackslash{}(-3\textbackslash{}) to \textbackslash{}(3\textbackslash{}). The graph consists of triangular peaks centered at \textbackslash{}(x = -2, 0, \textbackslash{}text\{ and \} 2\textbackslash{}).

- At \textbackslash{}(x = -2\textbackslash{}), and \textbackslash{}(x = 2\textbackslash{}), the triangles peak at 0.25 on the y-axis.
- At \textbackslash{}(x = 0\textbackslash{}), the triangle peaks at 0.50 on the y-axis.
- The lines forming the triangles create a symmetric pattern around the y-axis, with each triangle having a base width of about 2 units on the x-axis.

Overall, the plot is symmetrical with three evenly spaced triangular peaks.

\#\# Dataset metadata

Quantum Phenomena; Physical Model Grounding Reasoning, Numerical Reasoning

\paragraph{Recorded answer/context.}
C: C: \textbackslash{}( 5.0 \textbackslash{}times 10\textasciicircum{}\{-3\} \textbackslash{})

\paragraph{Figure/image path.}
/root/proposal\_for\_physic/science-mango/phyx\_mini\_run/phyx\_data/test\_image/502.png

\paragraph{Formalization target.}
create a compiling Lean file with sorry bodies at `PhyXMiniProblems/problem\_phyx\_mini\_0502.lean`. The Lean declarations must preserve the physical quantities, dimensions or dimensional roles, figure labels, governing-law hypotheses, and final relation expressed by this problem.
Use Mathlib/Physlib names found through LeanExplore where available. If a physics API is missing, introduce faithful local abstractions rather than scalar placeholder aliases.

\begin{theorem}[Physics formalization target]
\label{thm:physics:phyx_mini_0502:target}
\lean{PhyXMiniProblems.ProblemPhyXMini0502.problem_phyx_mini_0502}
\uses{def:physics:phyx-mini-0502:phyxminiproblems-problemphyxmini0502-electronpositionexperiment, def:physics:phyx-mini-0502:phyxminiproblems-problemphyxmini0502-satisfiesbornpositiondensitylaw, def:physics:phyx-mini-0502:phyxminiproblems-problemphyxmini0502-matchesproblemstatementandfigure, def:physics:phyx-mini-0502:phyxminiproblems-problemphyxmini0502-requestedcentralstrip, def:physics:phyx-mini-0502:phyxminiproblems-problemphyxmini0502-usesnarrowstripprobabilityapproximation, def:physics:phyx-mini-0502:phyxminiproblems-problemphyxmini0502-matchesanswerchoice}
The assigned autoformalize agent should translate this physics problem into Lean declarations in the covered file, with theorem and lemma proof bodies written as `by sorry`.
\end{theorem}
\begin{proof}
This is an autoformalization task, not a proof task. Produce statements that can later be proved by the physics prover without weakening the physical model.
\end{proof}

% --- Archon declaration topology begin ---
\section{Declaration topology}
\begin{definition}[Quantum Particle Kind]
  \label{def:physics:phyx-mini-0502:phyxminiproblems-problemphyxmini0502-quantumparticlekind}
  \lean{PhyXMiniProblems.ProblemPhyXMini0502.QuantumParticleKind}
  The particle species named in the scenario
\end{definition}
\begin{proof}
  This is the declared model definition of Quantum Particle Kind.
\end{proof}
\begin{definition}[Density Peak Label]
  \label{def:physics:phyx-mini-0502:phyxminiproblems-problemphyxmini0502-densitypeaklabel}
  \lean{PhyXMiniProblems.ProblemPhyXMini0502.DensityPeakLabel}
  Labels for the three triangular peaks visible in the graph
\end{definition}
\begin{proof}
  This is the declared model definition of Density Peak Label.
\end{proof}
\begin{definition}[Probability Density Curve Shape]
  \label{def:physics:phyx-mini-0502:phyxminiproblems-problemphyxmini0502-probabilitydensitycurveshape}
  \lean{PhyXMiniProblems.ProblemPhyXMini0502.ProbabilityDensityCurveShape}
  The qualitative curve shape read from the supplied graph
\end{definition}
\begin{proof}
  This is the declared model definition of Probability Density Curve Shape.
\end{proof}
\begin{definition}[Probability Density Figure]
  \label{def:physics:phyx-mini-0502:phyxminiproblems-problemphyxmini0502-probabilitydensityfigure}
  \lean{PhyXMiniProblems.ProblemPhyXMini0502.ProbabilityDensityFigure}
  \uses{def:physics:phyx-mini-0502:phyxminiproblems-problemphyxmini0502-densitypeaklabel, def:physics:phyx-mini-0502:phyxminiproblems-problemphyxmini0502-probabilitydensitycurveshape}
  The numerical and qualitative data printed in the probability-density graph. All positions, widths, and density values are named readouts in the units shown on the axes, rather than transparent aliases for physical quantities
\end{definition}
\begin{proof}
  This is the declared model definition of Probability Density Figure.
\end{proof}
\begin{definition}[Electron Position Experiment]
  \label{def:physics:phyx-mini-0502:phyxminiproblems-problemphyxmini0502-electronpositionexperiment}
  \lean{PhyXMiniProblems.ProblemPhyXMini0502.ElectronPositionExperiment}
  \uses{def:physics:phyx-mini-0502:phyxminiproblems-problemphyxmini0502-quantumparticlekind, def:physics:phyx-mini-0502:phyxminiproblems-problemphyxmini0502-probabilitydensityfigure}
  Physical state of the one-dimensional electron-position experiment. waveAmplitudePerSqrtMillimeter and densityPerMillimeter are calibrated coordinate representations of the wave amplitude and density. Physlib's QuantumMechanics.OneDimension.HilbertSpace.MemHS ensures that the amplitude is a square-integrable representative of a one-dimensional quantum state. The actual dimensionless position probability is the normalized measure positionProbability; it is not defined from an answer choice
\end{definition}
\begin{proof}
  This is the declared model definition of Electron Position Experiment.
\end{proof}
\begin{definition}[Satisfies Born Position Density Law]
  \label{def:physics:phyx-mini-0502:phyxminiproblems-problemphyxmini0502-satisfiesbornpositiondensitylaw}
  \lean{PhyXMiniProblems.ProblemPhyXMini0502.SatisfiesBornPositionDensityLaw}
  \uses{def:physics:phyx-mini-0502:phyxminiproblems-problemphyxmini0502-electronpositionexperiment}
  The Born position-density law and its measure-theoretic interpretation. Lebesgue volume is taken on the numerical millimetre coordinate, so a density readout has the displayed reciprocal-millimetre role. ENNReal.ofReal embeds the nonnegative real density supplied by the squared complex amplitude
\end{definition}
\begin{proof}
  This is the declared model definition of Satisfies Born Position Density Law.
\end{proof}
\begin{definition}[Matches Problem Statement And Figure]
  \label{def:physics:phyx-mini-0502:phyxminiproblems-problemphyxmini0502-matchesproblemstatementandfigure}
  \lean{PhyXMiniProblems.ProblemPhyXMini0502.MatchesProblemStatementAndFigure}
  \uses{def:physics:phyx-mini-0502:phyxminiproblems-problemphyxmini0502-densitypeaklabel, def:physics:phyx-mini-0502:phyxminiproblems-problemphyxmini0502-electronpositionexperiment}
  Scenario facts and primary-image readouts from 502.png. The three peak centers are -2, 0, and 2 mm; the corresponding peak densities are 0.25, 0.50, and 0.25 mm⁻¹; every triangular base is 2 mm wide; and the plotted axis spans -3 to 3 mm. These data contain no strip probability and no answer-choice value
\end{definition}
\begin{proof}
  This is the declared model definition of Matches Problem Statement And Figure.
\end{proof}
\begin{definition}[Detection Strip Readout]
  \label{def:physics:phyx-mini-0502:phyxminiproblems-problemphyxmini0502-detectionstripreadout}
  \lean{PhyXMiniProblems.ProblemPhyXMini0502.DetectionStripReadout}
  A detector strip specified by its center and full width, both in mm
\end{definition}
\begin{proof}
  This is the declared model definition of Detection Strip Readout.
\end{proof}
\begin{definition}[region]
  \label{def:physics:phyx-mini-0502:phyxminiproblems-problemphyxmini0502-detectionstripreadout-region}
  \lean{PhyXMiniProblems.ProblemPhyXMini0502.DetectionStripReadout.region}
  \uses{def:physics:phyx-mini-0502:phyxminiproblems-problemphyxmini0502-detectionstripreadout}
  The closed interval of millimetre-coordinate readings covered by a strip
\end{definition}
\begin{proof}
  This is the declared model definition of region.
\end{proof}
\begin{definition}[strip Probability]
  \label{def:physics:phyx-mini-0502:phyxminiproblems-problemphyxmini0502-stripprobability}
  \lean{PhyXMiniProblems.ProblemPhyXMini0502.stripProbability}
  \uses{def:physics:phyx-mini-0502:phyxminiproblems-problemphyxmini0502-electronpositionexperiment, def:physics:phyx-mini-0502:phyxminiproblems-problemphyxmini0502-detectionstripreadout}
  The dimensionless probability assigned by the model to a detector strip
\end{definition}
\begin{proof}
  This is the declared model definition of strip Probability.
\end{proof}
\begin{definition}[requested Central Strip]
  \label{def:physics:phyx-mini-0502:phyxminiproblems-problemphyxmini0502-requestedcentralstrip}
  \lean{PhyXMiniProblems.ProblemPhyXMini0502.requestedCentralStrip}
  \uses{def:physics:phyx-mini-0502:phyxminiproblems-problemphyxmini0502-detectionstripreadout}
  The 0.010-mm-wide strip centered at x = 0.000 mm
\end{definition}
\begin{proof}
  This is the declared model definition of requested Central Strip.
\end{proof}
\begin{definition}[Uses Narrow Strip Probability Approximation]
  \label{def:physics:phyx-mini-0502:phyxminiproblems-problemphyxmini0502-usesnarrowstripprobabilityapproximation}
  \lean{PhyXMiniProblems.ProblemPhyXMini0502.UsesNarrowStripProbabilityApproximation}
  \uses{def:physics:phyx-mini-0502:phyxminiproblems-problemphyxmini0502-electronpositionexperiment, def:physics:phyx-mini-0502:phyxminiproblems-problemphyxmini0502-detectionstripreadout, def:physics:phyx-mini-0502:phyxminiproblems-problemphyxmini0502-stripprobability}
  The textbook narrow-strip approximation used by the multiple-choice problem. The premise also records that the strip width is positive and smaller than the base width of the central peak. Its probability equation is a general local density-times-width relation and does not state the requested numeric value
\end{definition}
\begin{proof}
  This is the declared model definition of Uses Narrow Strip Probability Approximation.
\end{proof}
\begin{definition}[Answer Choice]
  \label{def:physics:phyx-mini-0502:phyxminiproblems-problemphyxmini0502-answerchoice}
  \lean{PhyXMiniProblems.ProblemPhyXMini0502.AnswerChoice}
  The four answer-choice labels printed with the problem
\end{definition}
\begin{proof}
  This is the declared model definition of Answer Choice.
\end{proof}
\begin{definition}[probability Value]
  \label{def:physics:phyx-mini-0502:phyxminiproblems-problemphyxmini0502-answerchoice-probabilityvalue}
  \lean{PhyXMiniProblems.ProblemPhyXMini0502.AnswerChoice.probabilityValue}
  \uses{def:physics:phyx-mini-0502:phyxminiproblems-problemphyxmini0502-answerchoice}
  Dimensionless probability value displayed beside an answer choice
\end{definition}
\begin{proof}
  This is the declared model definition of probability Value.
\end{proof}
\begin{definition}[recorded Answer Choice]
  \label{def:physics:phyx-mini-0502:phyxminiproblems-problemphyxmini0502-recordedanswerchoice}
  \lean{PhyXMiniProblems.ProblemPhyXMini0502.recordedAnswerChoice}
  \uses{def:physics:phyx-mini-0502:phyxminiproblems-problemphyxmini0502-answerchoice}
  Dataset metadata records choice C; this constant is never a premise
\end{definition}
\begin{proof}
  This is the declared model definition of recorded Answer Choice.
\end{proof}
\begin{definition}[Matches Answer Choice]
  \label{def:physics:phyx-mini-0502:phyxminiproblems-problemphyxmini0502-matchesanswerchoice}
  \lean{PhyXMiniProblems.ProblemPhyXMini0502.MatchesAnswerChoice}
  \uses{def:physics:phyx-mini-0502:phyxminiproblems-problemphyxmini0502-electronpositionexperiment, def:physics:phyx-mini-0502:phyxminiproblems-problemphyxmini0502-detectionstripreadout, def:physics:phyx-mini-0502:phyxminiproblems-problemphyxmini0502-stripprobability, def:physics:phyx-mini-0502:phyxminiproblems-problemphyxmini0502-answerchoice}
  A choice matches when its displayed value equals the modeled strip probability
\end{definition}
\begin{proof}
  This is the declared model definition of Matches Answer Choice.
\end{proof}
\begin{lemma}[central Strip Probability eq five thousandths]
  \label{lem:physics:phyx-mini-0502:phyxminiproblems-problemphyxmini0502-centralstripprobability-eq-five-thousandths}
  \lean{PhyXMiniProblems.ProblemPhyXMini0502.centralStripProbability_eq_five_thousandths}
  \uses{def:physics:phyx-mini-0502:phyxminiproblems-problemphyxmini0502-electronpositionexperiment, def:physics:phyx-mini-0502:phyxminiproblems-problemphyxmini0502-satisfiesbornpositiondensitylaw, def:physics:phyx-mini-0502:phyxminiproblems-problemphyxmini0502-matchesproblemstatementandfigure, def:physics:phyx-mini-0502:phyxminiproblems-problemphyxmini0502-stripprobability, def:physics:phyx-mini-0502:phyxminiproblems-problemphyxmini0502-requestedcentralstrip, def:physics:phyx-mini-0502:phyxminiproblems-problemphyxmini0502-usesnarrowstripprobabilityapproximation}
  For the requested central strip, the figure's 0.50 mm⁻¹ density and the 0.010 mm width give the dimensionless probability 5.0 × 10⁻³ under the narrow-strip approximation
\end{lemma}
\begin{proof}
  The conclusion follows from the governing relations and figure readouts named in the statement of central Strip Probability eq five thousandths.
\end{proof}
% --- Archon declaration topology end ---
% --- Archon physics formalization source end ---
